% Trang bìa theo mẫu của giảng viên (BÌA MẪU TIỂU LUẬN - T.HÙNG.docx)
\begin{titlepage}
\thispagestyle{empty}
\begin{tikzpicture}[remember picture,overlay]
  \draw[line width=2.2pt] ([shift={(1.5cm,-1.2cm)}]current page.north west) rectangle ([shift={(-1.2cm,1.2cm)}]current page.south east);
  \draw[line width=0.7pt] ([shift={(1.62cm,-1.32cm)}]current page.north west) rectangle ([shift={(-1.32cm,1.32cm)}]current page.south east);
\end{tikzpicture}
\centering

\includegraphics[width=0.75\textwidth]{image/siu-logo.png}

\vspace{1.2cm}
{\Large\bfseries KHOA KỸ THUẬT \& KHOA HỌC MÁY TÍNH}

\vspace{2.2cm}
{\fontsize{20}{24}\selectfont\bfseries TIỂU LUẬN CUỐI KỲ}\\[0.3cm]
{\fontsize{19}{23}\selectfont\bfseries MÔN KIẾN TRÚC MÁY TÍNH}

\vspace{2.2cm}
{\fontsize{20}{24}\selectfont\bfseries ĐỀ TÀI:}\\[0.3cm]
{\fontsize{15}{19}\selectfont\bfseries CÁCH TIẾP CẬN CHỊU LỖI CHỦ ĐỘNG CHO ĐIỆN TOÁN ĐÁM MÂY DỰA TRÊN HỆ MỜ TAKAGI-SUGENO VÀ THUẬT TOÁN SIMULATED ANNEALING\par}

\vspace{1.6cm}
\begin{flushleft}
\hspace{2.2cm}\begin{minipage}{0.8\textwidth}\large\bfseries
SINH VIÊN: NÔNG VĂN TÌNH\\[0.2cm]
MSSV: 8480101250028\\[0.2cm]
LỚP: 25MCS2\\[0.2cm]
GIẢNG VIÊN: GS.TS. TRẦN CÔNG HÙNG
\end{minipage}
\end{flushleft}

\vfill
{\Large\bfseries HK2 - NĂM HỌC 2025-2026}
\vspace{0.8cm}
\end{titlepage}
\clearpage
